\documentclass[11pt]{amsart}
\usepackage{calc,amssymb,amsmath,ulem,stmaryrd,fullpage}
\usepackage{enumerate}
\usepackage{alltt}
\RequirePackage[dvipsnames,usenames]{color}
\let\mffont=\sf
\normalem
%\input{kmacros3.sty}
\input{xy}
\xyoption{all}
\usepackage{hyperref}
\usepackage{graphicx}
\usepackage[all,cmtip]{xy}
\usepackage{verbatim}
\usepackage[left=0.95in,top=0.95in,right=0.95in,bottom=0.95in]{geometry}
\newcommand{\tauCohomology}{T}
\newcommand{\FCohomology}{S}


\theoremstyle{theorem}
%\newtheorem*{mainthm*}{Main Theorem}

\newcommand{\note}[1]{\marginpar{\sffamily\tiny #1}}

\def\todo#1{\textcolor{Mahogany}%
{\sffamily\footnotesize\newline{\color{Mahogany}\fbox{\parbox{\textwidth-15pt}{#1}}}\newline}}

\renewcommand{\O}{\mathcal O}
\newcommand{\ie}{{\itshape i.e.} }
\newcommand{\roundup}[1]{\lceil #1 \rceil}
\newcommand{\rounddown}[1]{\lfloor #1 \rfloor}
\renewcommand{\to}[1][]{\xrightarrow{\ #1\ }}
\newcommand{\onto}[1][]{\protect{\xrightarrow{\ #1\ }\hspace{-0.8em}\rightarrow}}
\newcommand{\into}[1][]{\lhook \joinrel \xrightarrow{\ #1\ }}
%\newcommand{DBDual}[1]{{\underline{\omega}_{#1}}}
\newcommand{\DBDual}[1]{{{\uuline \omega}{}^{\mydot}_{#1}}}
\newcommand{\Tt}{{\mathfrak{T}}}
%\newcommand{\bn}{{\mathfrak{n}} }
\newcommand{\sol}[1]{ \vskip 6pt{\color{blue} {\bf Solution:} #1} \vskip 6pt}

\newcommand{\bR}{{\mathbb{R}}}
\newcommand{\va}{{\vec{a}}}
\newcommand{\vb}{{\vec{b}}}
\newcommand{\vzero}{{\vec{0}}}
\newcommand{\vi}{{\vec{i}}}
\newcommand{\vj}{{\vec{j}}}
\newcommand{\vk}{{\vec{k}}}
\newcommand{\vh}{{\vec{h}}}
\newcommand{\vr}{{\vec{r}}}
\newcommand{\vu}{{\vec{u}}}
\newcommand{\vv}{{\vec{v}}}
\newcommand{\vw}{{\vec{w}}}
\newcommand{\vx}{{\vec{x}}}
\newcommand{\vy}{{\vec{y}}}
\newcommand{\vz}{{\vec{z}}}
\newcommand{\vT}{{\vec{T}}}
\newcommand{\vN}{{\vec{N}}}
\newcommand{\vF}{{\vec{F}}}
\newcommand{\vG}{{\vec{G}}}
\newcommand{\bF}{\mathbb{F}}
\newcommand{\bQ}{\mathbb{Q}}
\newcommand{\bZ}{\mathbb{Z}}
\newcommand{\bA}{\mathbb{A}}


\begin{document}
\title{Worksheet \#1 -- Math 6130\\Fall 2018}
\author{Due Friday, August 31st}

\maketitle
You may work in groups of up to 3.  Only one worksheet needs to be turned in per group.

Recall the following results that we started the class with.
\vskip 9pt
{\bf Theorem.} {\it Suppose that $\phi : X \to \bA^m$ is a regular map.  Then $\overline{\phi(X)} = Z(\ker \phi^*)$.}
\vskip 9pt
{\bf Theorem.} {\it A regular map $\phi : X \to Y$ is an isomorphism if and only if $\phi^* : k[Y] \to k[X]$ is an isomorphism.}
\vskip 9pt
\noindent
{\bf 1.} Show that a regular map $\phi : X \to Y$ is dominant if and only if $\phi^*$ injects.  
\vskip 200pt
{\bf 2.}  Show that $\phi : X \to Y$ induces an isomorphism of $X$ onto a closed subset of $Y$ if and only if $\phi^*$ surjects.  
\noindent
\newpage
\noindent
{\bf 3.}  Let $X = Z(y-x^2) \subseteq \bA^2$.  Show that $X$ is a variety isomorphic to $\bA^1$.   Then show that $Y = Z(xy - 1)$ is not isomorphic to $\bA^1$.  
\vskip 300pt
\noindent
{\bf 4.}  Let $X = Z(y^2 - x^3) \subseteq \bA^2$.  Consider the regular map $\phi : \bA^1 \to X$ defined by $\phi(t) = (t^3, t^2)$.  Show that $\phi$ is a bijection but not an isomorphism.
\newpage
\noindent
{\bf 5.}  Let $f \in k[x,y]$ be an irreducible quadratic polynomial.  If $\mathrm{char} k \neq 2$, show that either $Z(f)$ is isomorphic to $\bA^1$ or isomorphic to $Z(xy-1)$.
\vskip 300pt
\noindent
{\bf 6.}  Suppose that $X$ is an irreducible topological space.  If $U \subseteq X$ is a non-empty open set, show that $U$ is also irreducible.  Further show that if $U_1$ and $U_2$ are two non-emtpy sets, that $U_1 \cap U_2$ is non-empty.

\end{document}

